% CardioWatch — manuscript draft (Overleaf-ready).
% Compiles with pdfLaTeX on Overleaf as-is (standard article class).
% For PSB submission, port the body into the World Scientific ws-procs template.
\documentclass[11pt]{article}

\usepackage[margin=1in]{geometry}
\usepackage{graphicx}
\usepackage{booktabs}
\usepackage{amsmath}
\usepackage{siunitx}
\usepackage[hidelinks]{hyperref}
\usepackage{authblk}
\usepackage{caption}

\title{Device-Agnostic Detection of Atrial Fibrillation on Consumer
Wearables: A Pre-Registered Cross-Device Evaluation}

\author[1]{Urmi Shah}
\affil[1]{Department of Computer Engineering, San Jose State University}
\date{}

\begin{document}
\maketitle

\begin{abstract}
Atrial fibrillation (AFib) is a common heart-rhythm disorder and a major,
largely preventable cause of stroke; many cases go undiagnosed. Consumer
smartwatches that record a single-lead electrocardiogram (ECG) could help find
these cases earlier, but a model trained on one type of device often fails on
another. We study this cross-device problem directly. On a pre-registered,
leakage-safe, patient-grouped hold-out from a hospital 12-lead corpus
(CPSC~2018, $n{=}1{,}376$ records, Lead~I only), we compare a lightweight,
device-agnostic model that uses only heartbeat-timing features (RR~+~random
forest) against a deep convolutional--recurrent network (CNN-LSTM) trained on
the same data. The two are statistically indistinguishable
(AUC $0.935$ vs.\ $0.949$; DeLong $p{=}0.33$, McNemar $p{=}0.40$, Holm-corrected),
and the timing-based model transfers without retraining to ambulatory Holter
recordings (AUC $0.907$) and to real Apple~Watch exports (accuracy $0.907$, 95\%
CI $0.80$--$0.96$). A deep model additionally trained on wearable data
(AliveCor) reaches AUC $0.975$ but is not a same-data comparison. As an
exploratory extension, we test whether AFib onset can be predicted from
pre-onset sinus rhythm; using personalized timing and atrial-ectopy features
across two Holter cohorts (104 patients) we obtain only modest, sub-clinical
discrimination (AUC $0.66$). All headline numbers come from a single
pre-registered run and are released for reproduction. Our results argue that for
wearable AFib screening, simple device-agnostic features are competitive with
deep models and generalize better across hardware.
\end{abstract}

\section{Introduction}
Atrial fibrillation is the most common sustained arrhythmia and a leading cause
of ischemic stroke, yet a large share of cases are silent and undiagnosed.
Because the clinical benefit of finding AFib is timely anticoagulation that
prevents strokes, tools that detect it during everyday life are valuable.
Smartwatches that record a single-lead ECG (for example, Apple~Watch Series~4
and later) put such a tool in tens of millions of hands. The ECG signature of
AFib---absent P~waves and irregular beat-to-beat timing---is visible in a single
lead, which makes single-lead detection feasible in principle.

A practical obstacle stands in the way. Models trained on one recording device
often degrade sharply on another because they learn hardware-specific waveform
quirks rather than the underlying cardiac pattern. We refer to this as the
\emph{cross-device gap}. This paper asks a focused question: for single-lead
AFib detection that must run on consumer wearables, do we need deep waveform
models, or are simple device-agnostic features enough? We make three
contributions.
\begin{enumerate}
  \item A \textbf{pre-registered, leakage-safe controlled comparison} of a
  device-agnostic timing model against a deep model, trained on the same data
  and scored on the same held-out records, with paired significance testing.
  \item A \textbf{cross-device evaluation} of the timing model on three unseen
  hardware types---hospital 12-lead, ambulatory Holter, and Apple~Watch---without
  retraining.
  \item An \textbf{honest exploratory study} of AFib \emph{onset} prediction from
  pre-onset sinus rhythm, reported with its limitations rather than overstated.
\end{enumerate}
We deliberately scope the work to AFib detection and stroke prevention; it does
not address heart-attack (myocardial infarction) prediction, which is a
different condition with a different ECG signature.

\section{Related Work}
Single-lead AFib detection has been studied extensively since the PhysioNet/CinC
2017 challenge on short AliveCor recordings~\cite{clifford2017af}, and large
prospective work established that smartwatch screening can surface undiagnosed
AFib~\cite{perez2019appleheart}. Closest to our setting,
Bahrami~Rad~et~al.~\cite{bahramirad2024crowdsourced} combine open detectors and
report that timing-based features transfer across AliveCor and Apple~Watch
better than learned waveform features. Our study differs by running a
\emph{pre-registered, same-data} statistical comparison of a timing model
against a deep model and by adding a separate, clearly labeled onset-prediction
analysis. Onset prediction itself has a long history beginning with the PAF
Prediction Challenge~\cite{moody2001paf}; we revisit it on continuous Holter
data with patient-grouped evaluation.

\section{Data and Methods}

\subsection{Datasets}
We use five public or personally exported sources. CPSC~2018 (within the
PhysioNet/CinC~2020 corpus~\cite{alday2020classification}) provides hospital
12-lead ECGs, from which we take Lead~I only. PhysioNet/CinC~2017 provides
single-lead AliveCor recordings~\cite{clifford2017af}. The MIT-BIH Atrial
Fibrillation Database (\texttt{afdb}) and the Long-Term AF Database
(\texttt{ltafdb}) provide annotated ambulatory Holter
recordings~\cite{moody1983afdb,petrutiu2007ltaf,goldberger2000physionet}. A
small set of personal Apple~Watch exports (54 recordings, four volunteers, one
confirmed AFib) is used only for real-world validation; its labels are the
watch's own classifier output and are not independent clinical adjudication. A
separate tabular clinical dataset is used to train a clinical risk model that we
report as a distinct cohort.

\subsection{Models}
The \textbf{device-agnostic model} (RR~+~RF) detects AFib from heartbeat-timing
features only---coefficient of variation, RMSSD, pNN50, sample entropy, and
related statistics of the RR-interval series---fed to a random forest. Because
these features measure timing rather than waveform shape, they are insensitive
to device sampling rate and gain. The \textbf{deep model} (CNN-LSTM) is a
three-block 1D convolutional encoder followed by a two-layer LSTM and a linear
head, trained on 10-second Lead~I windows. A clinical risk model (random forest
and gradient boosting on tabular risk factors) is trained and reported
separately, since it uses a different patient population and is not fused into a
device-deployable score.

\subsection{Pre-registered evaluation protocol}
To avoid the common pitfall of comparing models trained on different data, we
fixed the analysis before training. We created a deterministic, label-stratified
hold-out of CPSC records (20\%, fixed seed) after checking for one-record-per-
patient, and excluded these records from \emph{both} models' training. The
device-agnostic model and a CPSC-only deep model were then scored on the
identical hold-out. Decision thresholds were frozen on training/validation data,
never on the hold-out. The primary discrimination test is the paired DeLong test
on AUC~\cite{delong1988comparing,sun2014fast}; the primary threshold-level test
is McNemar's test on correctness vectors. A single pre-registered Holm
family~\cite{holm1979simple} (DeLong plus fixed-threshold McNemar across the
model-versus-model pairs) controls multiple comparisons; matched-specificity and
random-baseline comparisons are reported as raw sensitivity and sanity checks.
Class imbalance during training was handled with SMOTE applied \emph{inside}
each cross-validation fold to prevent leakage~\cite{chawla2002smote}.
Proportion metrics carry Wilson confidence intervals.

\subsection{Exploratory onset prediction}
Separately, we test whether imminent AFib can be predicted from sinus rhythm.
From continuous Holter records we label 30-second sinus windows as
\emph{pre-onset} (a true onset follows within 30 minutes) or \emph{control} (no
onset within 60 minutes), excluding windows already in AFib. Features combine
RR/HRV statistics with neurokit2~\cite{makowski2021neurokit} atrial-ectopy and
expanded HRV markers, each z-scored against the patient's own control-window
baseline (personalization). Evaluation uses patient-grouped cross-validation
pooled over \texttt{afdb} and \texttt{ltafdb}; the operating threshold is set at
a fixed control specificity on out-of-fold scores, and we report onset
sensitivity, false alarms per hour, and the distribution of lead times.

\section{Results}

\subsection{Clinical risk model (separate cohort)}
Table~\ref{tab:clinical} reports the tabular clinical models on their own test
split. These are shown for completeness and are not combined with the ECG models.

\begin{table}[t]
\centering
\caption{Clinical risk models on the held-out clinical test split (distinct
cohort from the ECG data; not a cross-modal comparison).}
\label{tab:clinical}
\begin{tabular}{lcccc}
\toprule
Model & AUC & Recall & Specificity & F1 \\
\midrule
Random forest & 0.954 & 0.941 & 0.854 & 0.914 \\
Gradient boosting & 0.927 & 0.961 & 0.634 & 0.852 \\
\bottomrule
\end{tabular}
\end{table}

\subsection{Controlled ECG comparison}
On the pre-registered hold-out (Table~\ref{tab:ecg}), the device-agnostic timing
model and the CPSC-only deep model are statistically indistinguishable
(Table~\ref{tab:stats}): the AUC difference is $-0.013$ (95\% CI $-0.032$ to
$0.006$), DeLong $p{=}0.33$ and McNemar $p{=}0.40$ after Holm correction. The
deep model additionally trained on wearable data reaches AUC $0.975$ and is
significantly higher, but this is a deployment comparison with more and
different training data, not a same-data test. A random baseline confirms the
pipeline (AUC $0.502$). Figure~\ref{fig:roc} shows the ROC curves.

\begin{table}[t]
\centering
\caption{ECG models on the pre-registered CPSC hold-out ($n{=}1{,}376$, Lead~I).
Thresholds frozen on non-hold-out data.}
\label{tab:ecg}
\begin{tabular}{lccccc}
\toprule
Model & AUC & Recall & Specificity & F1 & Threshold \\
\midrule
RR~+~RF (device-agnostic) & 0.935 & 0.803 & 0.879 & 0.679 & 0.40 \\
CNN-LSTM (CPSC only) & 0.949 & 0.959 & 0.832 & 0.701 & 0.40 \\
CNN-LSTM (combined, deployment) & 0.975 & 0.918 & 0.941 & 0.837 & 0.40 \\
Random baseline & 0.502 & 0.508 & 0.511 & 0.269 & 0.50 \\
\bottomrule
\end{tabular}
\end{table}

\begin{table}[t]
\centering
\caption{Paired statistical comparisons on the hold-out. $\Delta$AUC and DeLong
test the discrimination difference; McNemar tests correctness at frozen
thresholds. $p$(Holm) is corrected across the model-versus-model family.}
\label{tab:stats}
\begin{tabular}{lcccc}
\toprule
Comparison & $\Delta$AUC (95\% CI) & DeLong $p$(Holm) & McNemar $p$(Holm) \\
\midrule
RR~+~RF vs.\ CNN-LSTM (CPSC) & $-0.013$ ($-0.032$, $0.006$) & 0.33 & 0.40 \\
RR~+~RF vs.\ CNN-LSTM (combined) & $-0.040$ ($-0.053$, $-0.027$) & $1.5\times10^{-8}$ & $1.7\times10^{-10}$ \\
CNN-LSTM (CPSC) vs.\ (combined) & $-0.027$ ($-0.040$, $-0.013$) & $2.8\times10^{-4}$ & $1.7\times10^{-17}$ \\
\bottomrule
\end{tabular}
\end{table}

\subsection{Cross-device generalization}
Without any retraining, the device-agnostic model transfers across hardware
(Table~\ref{tab:crossdevice}): AUC $0.907$ on ambulatory Holter (MIT-BIH,
28{,}104 windows from 25 patients) and accuracy $0.907$ (95\% CI $0.80$--$0.96$)
on real Apple~Watch exports. We stress that the Apple~Watch set has a single
AFib-positive recording, so its ``accuracy'' is dominated by specificity and its
sensitivity is effectively unmeasured; the labels are the watch's classifier
output, not clinical adjudication. These rows are therefore a real-world
plausibility check, not a sensitivity claim.

\begin{table}[t]
\centering
\caption{Zero-shot cross-device performance of the device-agnostic model.
Apple~Watch is exploratory (single positive; device-classifier labels).}
\label{tab:crossdevice}
\begin{tabular}{lcc}
\toprule
Test set & Device & Result \\
\midrule
CPSC hold-out & Hospital 12-lead (Lead~I) & AUC 0.935 \\
MIT-BIH AFib & Ambulatory Holter, 250\,Hz & AUC 0.907 \\
Apple~Watch (54 rec.) & Consumer wearable, 512\,Hz & Acc.\ 0.907 (0.80--0.96) \\
\bottomrule
\end{tabular}
\end{table}

\subsection{Detection timing}
On concatenated normal$\rightarrow$AFib recordings, the detector flags AFib
within the first analysis window (within 5--10 seconds of onset) in every
constructed transition, at a normal-rhythm false-alarm rate of about 8\% per
10-second window. This measures \emph{detection latency}, not advance warning;
because the onset is abrupt, no pre-onset signal exists to forecast.

\subsection{Exploratory onset prediction}
Predicting AFib before it begins is a harder, separate problem
(Table~\ref{tab:onset}, Figure~\ref{fig:onset}). Pooled across two Holter
cohorts (104 patients, patient-grouped cross-validation), personalized timing
and atrial-ectopy features give modest but consistent pre-onset discrimination
(AUC $0.66$; per-cohort $0.64$ and $0.66$). At about $0.5$ false alarms per hour
the model warns of $28\%$ of onsets a median of $1.2$ minutes ahead (IQR
$0.3$--$5.6$). Per-patient normalization was essential: without it, one cohort
sat at chance (AUC $0.54$). We report this as a real but sub-clinical signal,
not a clinical early-warning system.

\begin{table}[t]
\centering
\caption{Exploratory AFib onset prediction (pre-onset sinus; patient-grouped CV).}
\label{tab:onset}
\begin{tabular}{lccccc}
\toprule
Cohort & Patients & Window AUC & Onset sensitivity & Median lead (min) & FA/h \\
\midrule
Pooled & 104 & 0.663 & 0.28 (0.26--0.30) & 1.2 (IQR 0.3--5.6) & 0.48 \\
MIT-BIH (\texttt{afdb}) & 21 & 0.640 & 0.24 & 2.5 & 0.52 \\
Long-Term AF (\texttt{ltafdb}) & 83 & 0.662 & 0.29 & 1.2 & 0.48 \\
\bottomrule
\end{tabular}
\end{table}

\begin{figure}[t]
\centering
% Copy from docs/results/<run_id>/figures/roc_cpsc_holdout.png
\includegraphics[width=0.6\linewidth]{figures/roc_cpsc_holdout.png}
\caption{ROC curves on the pre-registered CPSC hold-out. The device-agnostic
model and the CPSC-only deep model overlap; the random baseline sits on the
diagonal.}
\label{fig:roc}
\end{figure}

\begin{figure}[t]
\centering
% Copy from docs/results/<run_id>/figures/onset_prediction_nk_lead_time.png
\includegraphics[width=0.6\linewidth]{figures/onset_prediction_nk_lead_time.png}
\caption{Distribution of lead times for warned AFib onsets in the exploratory
onset analysis.}
\label{fig:onset}
\end{figure}

\section{Discussion}
The central result is a negative one in the useful sense: on matched data, a
simple timing-based model is not measurably worse than a deep waveform model for
single-lead AFib detection, and it transfers across hardware without retraining.
For a wearable screening setting---where the device, sampling rate, and noise
profile vary---this favors the simpler, interpretable model. The combined deep
model scores higher, but only with extra wearable training data; that is a
statement about data, not architecture.

Several limitations bound our claims. The clinical and ECG models use different
populations and are not a cross-modal ablation. The Apple~Watch validation is
small, has one positive case, and relies on device-generated labels, so it
supports plausibility rather than sensitivity. Detection latency reflects
near-instant detection after an abrupt onset, not advance warning. The onset
analysis uses cohorts of selected AFib patients, shows only modest
discrimination, and is therefore framed as exploratory; richer atrial markers
(such as P-wave morphology), longer pre-onset context, and purpose-built data
are the natural next steps.

\section{Conclusion}
For single-lead AFib detection on consumer wearables, device-agnostic timing
features match a deep model on matched data and generalize better across
devices, making them a strong, interpretable default for stroke-prevention
screening. Predicting AFib before it starts remains hard: our personalized
features yield a real but sub-clinical early-warning signal that motivates
future work rather than a deployable alarm.

\section*{Reproducibility}
All reported numbers come from one pre-registered run. Code, the analysis plan,
the hold-out manifest checksum, and the compact result files are available in the
project repository; a versioned snapshot will be archived with a DOI.

\bibliographystyle{unsrt}
\bibliography{references}

\end{document}
